% !TEX root = ../Main.tex
\chapter{静止开始的匀加速直线运动的"比"}

物体从静止出发、做匀加速直线运动，这是高一动力学最基本的模型。它的核心就是一张 $v$-$t$ 图像：一条从原点出发的斜直线，斜率是加速度 $a$，图线下的面积就是位移。所有"比例"都能从这张图直接读出来，不必死记。

\section{核心：$v$-$t$ 图像}

设物体从静止出发、加速度恒为 $a$。任意时刻速度 $v=at$，到 $t$ 时刻的位移就是图线下三角形的面积
\[
s=\tfrac{1}{2}at^2.
\]

\begin{center}
\begin{tikzpicture}[>=Stealth, scale=1]
    \draw[->] (-0.3,0) -- (5.5,0) node[right] {$t$};
    \draw[->] (0,-0.3) -- (0,3.4) node[above] {$v$};
    \draw[thick, blue] (0,0) -- (5,3);
    \foreach \x/\h/\l in {1/0.6/T, 2/1.2/2T, 3/1.8/3T, 4/2.4/4T} {
        \draw[dashed] (\x,0) -- (\x,\h);
        \node[below] at (\x,-0.05) {\small $\l$};
    }
    \draw[dashed] (5,0) -- (5,3);
    \node[below] at (5,-0.05) {\small $nT$};
\end{tikzpicture}
\end{center}

把时间轴按相等的间隔 $T$ 分段（$T,2T,3T,\dots,nT$），就有三组常用的比例。

\section{三组比例}

\state{三组比例（从静止出发的匀加速）}{
    把时间按 $T$ 等分，记：

    \begin{itemize}
        \item $v_k$：第 $k$ 个时刻末（即 $t=kT$）的瞬时速度；
        \item $s_k$：从出发到 $t=kT$ 的\textbf{累计位移}；
        \item $s_{(k-1)k}$：第 $k$ 段时间（即第 $k$ 个 $T$）内的位移。
    \end{itemize}

    \[
    \begin{aligned}
    v_1:v_2:v_3:\cdots:v_n &= \rule{4cm}{0.4pt}\\[6pt]
    s_1:s_2:s_3:\cdots:s_n &= \rule{4cm}{0.4pt}\\[6pt]
    s_{12}:s_{23}:s_{34}:\cdots:s_{(n-1)n} &= \rule{4cm}{0.4pt}
    \end{aligned}
    \]
}
